% Author: Izaak Neutelings (October 2020)
% Instructions: To compile, include data files from
%   https://github.com/IzaakWN/CodeSnippets/tree/master/LaTeX/TikZ/physics/dynamics_pendulum
% Other sources:
%   https://www.scielo.br/scielo.php?script=sci_arttext&pid=S1806-11172007000400024
%   https://www.researchgate.net/publication/262746594_Exact_solution_for_the_nonlinear_pendulum
%   https://tex.stackexchange.com/questions/545590/phase-portrait-of-van-der-pol-oscillator-with-pgfplots
%   http://matlab.cheme.cmu.edu/2011/08/09/phase-portraits-of-a-system-of-odes/
\documentclass[border=3pt,tikz]{standalone}
\usepackage{physics}
\usepackage{siunitx}
\sisetup{detect-all} % to allow \pi in \SI
\usepackage{tikz,pgfplots}
\usepackage[outline]{contour} % glow around text
\usetikzlibrary{calc}
\usetikzlibrary{angles,quotes} % for pic
\usetikzlibrary{arrows.meta}
\tikzset{>=latex} % for LaTeX arrow head
\contourlength{1.2pt}

\colorlet{xcol}{blue!70!black}
\colorlet{vcol}{green!60!black}
\colorlet{myred}{red!70!black}
\colorlet{myblue}{blue!70!black}
\colorlet{mygreen}{green!70!black}
\colorlet{mydarkred}{myred!70!black}
\colorlet{mydarkblue}{myblue!60!black}
\colorlet{mydarkgreen}{mygreen!60!black}
\colorlet{acol}{red!50!blue!80!black!80}
\tikzstyle{CM}=[red!40!black,fill=red!80!black!80]
\tikzstyle{xline}=[xcol,thick,smooth]
\tikzstyle{mass}=[line width=0.6,red!30!black,fill=red!40!black!10,rounded corners=1,
                  top color=red!40!black!20,bottom color=red!40!black!10,shading angle=20]
\tikzstyle{faded mass}=[dashed,line width=0.1,red!30!black!40,fill=red!40!black!10,rounded corners=1,
                        top color=red!40!black!10,bottom color=red!40!black!10,shading angle=20]
\tikzstyle{rope}=[brown!70!black,very thick,line cap=round]
\def\rope#1{ \draw[black,line width=1.4] #1; \draw[rope,line width=1.1] #1; }
\tikzstyle{force}=[->,myred,very thick,line cap=round]
\tikzstyle{velocity}=[->,vcol,very thick,line cap=round]
\tikzstyle{Fproj}=[force,myred!40]
\tikzstyle{myarr}=[-{Latex[length=3,width=2]},thin]
\def\tick#1#2{\draw[thick] (#1)++(#2:0.12) --++ (#2-180:0.24)}
\DeclareMathOperator{\sn}{sn}
\DeclareMathOperator{\cn}{cn}
\DeclareMathOperator{\dn}{dn}
\def\N{80} % number of samples in plots


\begin{document}
% PENDULUM
\def\L{2.8}  % string length
\def\ang{28} % angle string
\def\R{0.25} % ball radius
\def\F{1.0}  % force magnitude
% PENDULUM FORCES
\begin{tikzpicture}
    \message{^^JPendulum forces}
    \def\F{1.4}  % force magnitude
    \coordinate (O) at (0,0);
    \coordinate (FT) at (90+\ang:{\F*cos(\ang)});
    \coordinate (FG) at (-90:\F);
    \coordinate (FGx) at (-90+\ang:{\F*cos(\ang)});
    \coordinate (FGy) at (180+\ang:{\F*sin(\ang)});
    \draw[dashed,myred!60!black] (FGy) -- (FG); %-- (FGx);
    \draw[force] (O) -- (FT) node[midway,above right=-2] {$\vb{T}$};
    \draw[Fproj] (O) -- (FGy) node[pos=0.4,above left=-3,scale=0.86] {$mg\sin\theta$};
    \draw[Fproj] (O) -- (FGx) node[pos=0.5,above right=-3,scale=0.86] {$mg\cos\theta$};
    \draw[force] (O) -- (FG) node[right=0] {$m\vb{g}$};
    \draw pic[myarr,"$\theta$",xcol,draw=xcol,angle radius=14,angle eccentricity=1.45] {angle=FG--O--FGx};
\end{tikzpicture}


\end{document}